% ============================================================================
% SLIDES - Physik Klasse 6: Weg-Zeit-Diagramme
% ============================================================================
% Struktur pro Thema:
%   1. Übung (nur Frage)
%   2. Übung + Lösung
%   3. Übung + Lösung + AB-Verweis
%   4. ML-Folie (NUR AB-Lösung)
% Arbeitspfad: Kasten 1 → 2 → 3 → Aufgabe 1 → 2 → 3 → 4
% ============================================================================

\documentclass[aspectratio=169,11pt]{beamer}

\usepackage[utf8]{inputenc}
\usepackage[T1]{fontenc}
\usepackage[ngerman]{babel}
\usepackage{tikz}
\usepackage{tcolorbox}
\usepackage{amssymb}

\usetheme{default}
\usecolortheme{default}

\definecolor{physik}{HTML}{2c5aa0}
\definecolor{gridcolor}{HTML}{c0d0e8}
\definecolor{hinweis}{HTML}{b35c00}
\definecolor{loesung}{HTML}{c62828}

\setbeamercolor{frametitle}{fg=physik}
\setbeamercolor{title}{fg=physik}
\setbeamercolor{structure}{fg=physik}

\setbeamertemplate{navigation symbols}{}
\setbeamertemplate{footline}{
    \hfill\textcolor{gray}{\footnotesize\insertframenumber/\inserttotalframenumber}\hspace{1em}\vspace{0.5em}
}

\newtcolorbox{abhinweis}{
    colback=hinweis!10, colframe=hinweis,
    fonttitle=\bfseries\small, title={\textcolor{white}{Arbeitsblatt}},
    boxrule=1.5pt, arc=2pt, left=6pt, right=6pt, top=4pt, bottom=4pt,
    coltitle=white, colbacktitle=hinweis
}

\newtcolorbox{mlbox}{
    colback=loesung!5, colframe=loesung,
    fonttitle=\bfseries\small, title={\textcolor{white}{Musterlösung}},
    boxrule=2pt, arc=2pt, left=8pt, right=8pt, top=6pt, bottom=6pt,
    coltitle=white, colbacktitle=loesung
}

\newcommand{\luecke}[1]{\underline{\hspace{#1}}}

\title{\textbf{Weg-Zeit-Diagramme}}
\subtitle{Bewegungen sichtbar machen}
\author{Physik Klasse 6}
\date{}

\begin{document}

% === TITEL ===
\begin{frame}
    \titlepage
\end{frame}

% ============================================================================
% KASTEN 1: Achsen
% ============================================================================

% --- Kasten 1: Frage ---
\begin{frame}{Was gehört auf welche Achse?}
    \begin{columns}
        \begin{column}{0.5\textwidth}
            Wir wollen Bewegungen als \textbf{Diagramm} darstellen.

            \vspace{0.8em}
            \textbf{Messwerte:}
            \begin{tabular}{|c|c|c|c|c|c|}
                \hline
                $t$ in s & 0 & 2 & 4 & 6 & 8 \\
                \hline
                $s$ in m & 0 & 10 & 20 & 30 & 40 \\
                \hline
            \end{tabular}

            \vspace{1em}
            \textbf{Frage:}\\
            Was kommt auf die waagerechte Achse?\\
            Was kommt auf die senkrechte Achse?
        \end{column}
        \begin{column}{0.45\textwidth}
            \begin{tikzpicture}[scale=0.55]
                \draw[gridcolor, very thin] (0,0) grid (8,5);
                \draw[->, very thick, physik] (0,0) -- (8.5,0) node[right] {\textbf{?}};
                \draw[->, very thick, physik] (0,0) -- (0,5.5) node[above] {\textbf{?}};
            \end{tikzpicture}
        \end{column}
    \end{columns}
\end{frame}

% --- Kasten 1: Lösung ---
\begin{frame}{Was gehört auf welche Achse?}
    \begin{columns}
        \begin{column}{0.5\textwidth}
            Wir wollen Bewegungen als \textbf{Diagramm} darstellen.

            \vspace{0.8em}
            \textbf{Messwerte:}
            \begin{tabular}{|c|c|c|c|c|c|}
                \hline
                $t$ in s & 0 & 2 & 4 & 6 & 8 \\
                \hline
                $s$ in m & 0 & 10 & 20 & 30 & 40 \\
                \hline
            \end{tabular}

            \vspace{1em}
            \textcolor{loesung}{\textbf{Lösung:}\\
            Waagerecht: \textbf{Zeit} $t$ in s\\
            Senkrecht: \textbf{Weg} $s$ in m}
        \end{column}
        \begin{column}{0.45\textwidth}
            \begin{tikzpicture}[scale=0.55]
                \draw[gridcolor, very thin] (0,0) grid (8,5);
                \draw[->, very thick, loesung] (0,0) -- (8.5,0) node[right] {\textbf{$t$ in s}};
                \draw[->, very thick, loesung] (0,0) -- (0,5.5) node[above] {\textbf{$s$ in m}};
                \foreach \x/\y in {0/0, 2/1, 4/2, 6/3, 8/4}
                    \filldraw[loesung] (\x,\y) circle (3pt);
                \draw[loesung, very thick] (0,0) -- (8,4);
            \end{tikzpicture}
        \end{column}
    \end{columns}
\end{frame}

% --- Kasten 1: + AB-Verweis ---
\begin{frame}{Was gehört auf welche Achse?}
    \begin{columns}
        \begin{column}{0.5\textwidth}
            Wir wollen Bewegungen als \textbf{Diagramm} darstellen.

            \vspace{0.8em}
            \textbf{Messwerte:}
            \begin{tabular}{|c|c|c|c|c|c|}
                \hline
                $t$ in s & 0 & 2 & 4 & 6 & 8 \\
                \hline
                $s$ in m & 0 & 10 & 20 & 30 & 40 \\
                \hline
            \end{tabular}

            \vspace{1em}
            \textcolor{loesung}{\textbf{Lösung:}\\
            Waagerecht: \textbf{Zeit} $t$ in s\\
            Senkrecht: \textbf{Weg} $s$ in m}
        \end{column}
        \begin{column}{0.45\textwidth}
            \begin{tikzpicture}[scale=0.55]
                \draw[gridcolor, very thin] (0,0) grid (8,5);
                \draw[->, very thick, loesung] (0,0) -- (8.5,0) node[right] {\textbf{$t$ in s}};
                \draw[->, very thick, loesung] (0,0) -- (0,5.5) node[above] {\textbf{$s$ in m}};
                \foreach \x/\y in {0/0, 2/1, 4/2, 6/3, 8/4}
                    \filldraw[loesung] (\x,\y) circle (3pt);
                \draw[loesung, very thick] (0,0) -- (8,4);
            \end{tikzpicture}
        \end{column}
    \end{columns}

    \vspace{0.3em}
    \begin{abhinweis}
        \textbf{Kasten 1:} Ergänze die Achsenbeschriftung
    \end{abhinweis}
\end{frame}

% --- ML Kasten 1 ---
\begin{frame}{Musterlösung: Kasten 1}
    \begin{center}
        \begin{mlbox}
            \textbf{Kasten 1: Aufbau eines Weg-Zeit-Diagramms}

            \vspace{0.5em}
            \begin{itemize}
                \item Die \textbf{waagerechte Achse} (x-Achse) zeigt die \textcolor{loesung}{\textbf{Zeit}} $t$ in \textcolor{loesung}{\textbf{s}}.
                \item Die \textbf{senkrechte Achse} (y-Achse) zeigt den \textcolor{loesung}{\textbf{Weg}} $s$ in \textcolor{loesung}{\textbf{m}}.
            \end{itemize}
        \end{mlbox}
    \end{center}
\end{frame}

% ============================================================================
% KASTEN 2: Steigung
% ============================================================================

% --- Kasten 2: Frage ---
\begin{frame}{Wer ist schneller?}
    \begin{columns}
        \begin{column}{0.5\textwidth}
            Zwei Personen bewegen sich.

            \vspace{1em}
            \textbf{Schau auf die Geraden:}\\
            Welche gehört zum schnelleren Objekt?

            \vspace{1em}
            \textbf{Überlege:}
            \begin{itemize}
                \item Was bedeutet \textbf{steil}?
                \item Was bedeutet \textbf{flach}?
            \end{itemize}
        \end{column}
        \begin{column}{0.45\textwidth}
            \begin{tikzpicture}[scale=0.5]
                \draw[gridcolor, very thin] (0,0) grid (8,5);
                \draw[->, thick] (0,0) -- (8.5,0) node[right] {$t$ in s};
                \draw[->, thick] (0,0) -- (0,5.5) node[above] {$s$ in m};
                \foreach \x in {0,2,4,6,8} \draw (\x,0) node[below] {\small \x};
                \foreach \y/\l in {0/0,1/10,2/20,3/30,4/40} \draw (0,\y) node[left] {\small \l};
                \draw[red!70!black, very thick] (0,0) -- (5.3,4);
                \node[red!70!black, above] at (3,3.5) {\small Radfahrer};
                \draw[physik, very thick] (0,0) -- (8,2);
                \node[physik, right] at (8.2,2) {\small Fußgänger};
            \end{tikzpicture}
        \end{column}
    \end{columns}
\end{frame}

% --- Kasten 2: Lösung ---
\begin{frame}{Wer ist schneller?}
    \begin{columns}
        \begin{column}{0.5\textwidth}
            Zwei Personen bewegen sich.

            \vspace{1em}
            \textcolor{loesung}{\textbf{Lösung:}\\
            Der \textbf{Radfahrer} ist schneller!\\
            Seine Gerade ist \textbf{steiler}.}

            \vspace{0.8em}
            \textcolor{loesung}{\textbf{Faustregel:}\\
            Steil = schnell\\
            Flach = langsam}
        \end{column}
        \begin{column}{0.45\textwidth}
            \begin{tikzpicture}[scale=0.5]
                \draw[gridcolor, very thin] (0,0) grid (8,5);
                \draw[->, thick] (0,0) -- (8.5,0) node[right] {$t$ in s};
                \draw[->, thick] (0,0) -- (0,5.5) node[above] {$s$ in m};
                \foreach \x in {0,2,4,6,8} \draw (\x,0) node[below] {\small \x};
                \foreach \y/\l in {0/0,1/10,2/20,3/30,4/40} \draw (0,\y) node[left] {\small \l};
                \draw[loesung, very thick] (0,0) -- (5.3,4);
                \node[loesung, above] at (3,3.5) {\small schnell};
                \draw[physik, very thick] (0,0) -- (8,2);
                \node[physik, right] at (8.2,2) {\small langsam};
            \end{tikzpicture}
        \end{column}
    \end{columns}
\end{frame}

% --- Kasten 2: + AB-Verweis ---
\begin{frame}{Wer ist schneller?}
    \begin{columns}
        \begin{column}{0.5\textwidth}
            Zwei Personen bewegen sich.

            \vspace{1em}
            \textcolor{loesung}{\textbf{Lösung:}\\
            Der \textbf{Radfahrer} ist schneller!\\
            Seine Gerade ist \textbf{steiler}.}

            \vspace{0.8em}
            \textcolor{loesung}{\textbf{Faustregel:}\\
            Steil = schnell\\
            Flach = langsam}
        \end{column}
        \begin{column}{0.45\textwidth}
            \begin{tikzpicture}[scale=0.5]
                \draw[gridcolor, very thin] (0,0) grid (8,5);
                \draw[->, thick] (0,0) -- (8.5,0) node[right] {$t$ in s};
                \draw[->, thick] (0,0) -- (0,5.5) node[above] {$s$ in m};
                \foreach \x in {0,2,4,6,8} \draw (\x,0) node[below] {\small \x};
                \foreach \y/\l in {0/0,1/10,2/20,3/30,4/40} \draw (0,\y) node[left] {\small \l};
                \draw[loesung, very thick] (0,0) -- (5.3,4);
                \node[loesung, above] at (3,3.5) {\small schnell};
                \draw[physik, very thick] (0,0) -- (8,2);
                \node[physik, right] at (8.2,2) {\small langsam};
            \end{tikzpicture}
        \end{column}
    \end{columns}

    \vspace{0.3em}
    \begin{abhinweis}
        \textbf{Kasten 2:} Ergänze: steil = ?, flach = ?
    \end{abhinweis}
\end{frame}

% --- ML Kasten 2 ---
\begin{frame}{Musterlösung: Kasten 2}
    \begin{center}
        \begin{mlbox}
            \textbf{Kasten 2: Schnell oder langsam?}

            \vspace{0.5em}
            Die \textbf{Steigung} der Geraden zeigt die Geschwindigkeit:
            \begin{itemize}
                \item \textbf{Steile} Gerade = \textcolor{loesung}{\textbf{schnelle}} Bewegung
                \item \textbf{Flache} Gerade = \textcolor{loesung}{\textbf{langsame}} Bewegung
            \end{itemize}
        \end{mlbox}
    \end{center}
\end{frame}

% ============================================================================
% KASTEN 3: Vier Bewegungsarten
% ============================================================================

% --- Kasten 3: Frage ---
\begin{frame}{Vier verschiedene Linien -- Was bedeuten sie?}
    \begin{center}
        \textbf{Ordne zu: schnell, langsam, Stillstand, ungleichförmig}

        \vspace{0.5em}
        \begin{tikzpicture}[scale=0.42]
            \begin{scope}[shift={(0,0)}]
                \draw[gridcolor, very thin] (0,0) grid (5,4);
                \draw[->, thick] (0,0) -- (5.5,0); \draw[->, thick] (0,0) -- (0,4.5);
                \draw[physik, very thick] (0,0) -- (5,4);
                \node[below] at (2.5,-0.8) {\small \textbf{steil}};
                \node[below] at (2.5,-1.6) {\small = ?};
            \end{scope}
            \begin{scope}[shift={(7,0)}]
                \draw[gridcolor, very thin] (0,0) grid (5,4);
                \draw[->, thick] (0,0) -- (5.5,0); \draw[->, thick] (0,0) -- (0,4.5);
                \draw[physik, very thick] (0,0) -- (5,1.5);
                \node[below] at (2.5,-0.8) {\small \textbf{flach}};
                \node[below] at (2.5,-1.6) {\small = ?};
            \end{scope}
            \begin{scope}[shift={(14,0)}]
                \draw[gridcolor, very thin] (0,0) grid (5,4);
                \draw[->, thick] (0,0) -- (5.5,0); \draw[->, thick] (0,0) -- (0,4.5);
                \draw[physik, very thick] (0,2) -- (5,2);
                \node[below] at (2.5,-0.8) {\small \textbf{waagerecht}};
                \node[below] at (2.5,-1.6) {\small = ?};
            \end{scope}
            \begin{scope}[shift={(21,0)}]
                \draw[gridcolor, very thin] (0,0) grid (5,4);
                \draw[->, thick] (0,0) -- (5.5,0); \draw[->, thick] (0,0) -- (0,4.5);
                \draw[physik, very thick] (0,0) .. controls (2,0.5) and (3,2) .. (5,4);
                \node[below] at (2.5,-0.8) {\small \textbf{Kurve}};
                \node[below] at (2.5,-1.6) {\small = ?};
            \end{scope}
        \end{tikzpicture}
    \end{center}
\end{frame}

% --- Kasten 3: Lösung ---
\begin{frame}{Vier verschiedene Linien -- Was bedeuten sie?}
    \begin{center}
        \textbf{Ordne zu: schnell, langsam, Stillstand, ungleichförmig}

        \vspace{0.5em}
        \begin{tikzpicture}[scale=0.42]
            \begin{scope}[shift={(0,0)}]
                \draw[gridcolor, very thin] (0,0) grid (5,4);
                \draw[->, thick] (0,0) -- (5.5,0); \draw[->, thick] (0,0) -- (0,4.5);
                \draw[physik, very thick] (0,0) -- (5,4);
                \node[below] at (2.5,-0.8) {\small \textbf{steil}};
                \node[below, loesung] at (2.5,-1.6) {\small = \textbf{schnell}};
            \end{scope}
            \begin{scope}[shift={(7,0)}]
                \draw[gridcolor, very thin] (0,0) grid (5,4);
                \draw[->, thick] (0,0) -- (5.5,0); \draw[->, thick] (0,0) -- (0,4.5);
                \draw[physik, very thick] (0,0) -- (5,1.5);
                \node[below] at (2.5,-0.8) {\small \textbf{flach}};
                \node[below, loesung] at (2.5,-1.6) {\small = \textbf{langsam}};
            \end{scope}
            \begin{scope}[shift={(14,0)}]
                \draw[gridcolor, very thin] (0,0) grid (5,4);
                \draw[->, thick] (0,0) -- (5.5,0); \draw[->, thick] (0,0) -- (0,4.5);
                \draw[physik, very thick] (0,2) -- (5,2);
                \node[below] at (2.5,-0.8) {\small \textbf{waagerecht}};
                \node[below, loesung] at (2.5,-1.6) {\small = \textbf{Stillstand}};
            \end{scope}
            \begin{scope}[shift={(21,0)}]
                \draw[gridcolor, very thin] (0,0) grid (5,4);
                \draw[->, thick] (0,0) -- (5.5,0); \draw[->, thick] (0,0) -- (0,4.5);
                \draw[physik, very thick] (0,0) .. controls (2,0.5) and (3,2) .. (5,4);
                \node[below] at (2.5,-0.8) {\small \textbf{Kurve}};
                \node[below, loesung] at (2.5,-1.6) {\small = \textbf{ungleichf.}};
            \end{scope}
        \end{tikzpicture}
    \end{center}
\end{frame}

% --- Kasten 3: + AB-Verweis ---
\begin{frame}{Vier verschiedene Linien -- Was bedeuten sie?}
    \begin{center}
        \begin{tikzpicture}[scale=0.38]
            \begin{scope}[shift={(0,0)}]
                \draw[gridcolor, very thin] (0,0) grid (5,4);
                \draw[->, thick] (0,0) -- (5.5,0); \draw[->, thick] (0,0) -- (0,4.5);
                \draw[physik, very thick] (0,0) -- (5,4);
                \node[below] at (2.5,-0.7) {\small \textbf{steil}};
                \node[below, loesung] at (2.5,-1.5) {\small = \textbf{schnell}};
            \end{scope}
            \begin{scope}[shift={(7,0)}]
                \draw[gridcolor, very thin] (0,0) grid (5,4);
                \draw[->, thick] (0,0) -- (5.5,0); \draw[->, thick] (0,0) -- (0,4.5);
                \draw[physik, very thick] (0,0) -- (5,1.5);
                \node[below] at (2.5,-0.7) {\small \textbf{flach}};
                \node[below, loesung] at (2.5,-1.5) {\small = \textbf{langsam}};
            \end{scope}
            \begin{scope}[shift={(14,0)}]
                \draw[gridcolor, very thin] (0,0) grid (5,4);
                \draw[->, thick] (0,0) -- (5.5,0); \draw[->, thick] (0,0) -- (0,4.5);
                \draw[physik, very thick] (0,2) -- (5,2);
                \node[below] at (2.5,-0.7) {\small \textbf{waagerecht}};
                \node[below, loesung] at (2.5,-1.5) {\small = \textbf{Stillstand}};
            \end{scope}
            \begin{scope}[shift={(21,0)}]
                \draw[gridcolor, very thin] (0,0) grid (5,4);
                \draw[->, thick] (0,0) -- (5.5,0); \draw[->, thick] (0,0) -- (0,4.5);
                \draw[physik, very thick] (0,0) .. controls (2,0.5) and (3,2) .. (5,4);
                \node[below] at (2.5,-0.7) {\small \textbf{Kurve}};
                \node[below, loesung] at (2.5,-1.5) {\small = \textbf{ungleichf.}};
            \end{scope}
        \end{tikzpicture}
    \end{center}

    \vspace{0.3em}
    \begin{abhinweis}
        \textbf{Kasten 3:} Ergänze die vier Bewegungsarten
    \end{abhinweis}
\end{frame}

% --- ML Kasten 3 ---
\begin{frame}{Musterlösung: Kasten 3}
    \begin{center}
        \begin{mlbox}
            \textbf{Kasten 3: Vier Bewegungsarten im Diagramm}

            \vspace{0.3em}
            \begin{tikzpicture}[scale=0.32]
                \begin{scope}[shift={(0,0)}]
                    \draw[gridcolor] (0,0) grid (4,3);
                    \draw[->, thick] (0,0) -- (4.3,0); \draw[->, thick] (0,0) -- (0,3.3);
                    \draw[physik, very thick] (0,0) -- (4,3);
                    \node[below] at (2,-0.6) {\tiny steil};
                    \node[below, loesung] at (2,-1.4) {\tiny = \textbf{schnell}};
                \end{scope}
                \begin{scope}[shift={(6,0)}]
                    \draw[gridcolor] (0,0) grid (4,3);
                    \draw[->, thick] (0,0) -- (4.3,0); \draw[->, thick] (0,0) -- (0,3.3);
                    \draw[physik, very thick] (0,0) -- (4,1.2);
                    \node[below] at (2,-0.6) {\tiny flach};
                    \node[below, loesung] at (2,-1.4) {\tiny = \textbf{langsam}};
                \end{scope}
                \begin{scope}[shift={(12,0)}]
                    \draw[gridcolor] (0,0) grid (4,3);
                    \draw[->, thick] (0,0) -- (4.3,0); \draw[->, thick] (0,0) -- (0,3.3);
                    \draw[physik, very thick] (0,1.5) -- (4,1.5);
                    \node[below] at (2,-0.6) {\tiny waager.};
                    \node[below, loesung] at (2,-1.4) {\tiny = \textbf{Stillstand}};
                \end{scope}
                \begin{scope}[shift={(18,0)}]
                    \draw[gridcolor] (0,0) grid (4,3);
                    \draw[->, thick] (0,0) -- (4.3,0); \draw[->, thick] (0,0) -- (0,3.3);
                    \draw[physik, very thick] (0,0) .. controls (1.5,0.5) and (2.5,2) .. (4,3);
                    \node[below] at (2,-0.6) {\tiny Kurve};
                    \node[below, loesung] at (2,-1.4) {\tiny = \textbf{ungleichf.}};
                \end{scope}
            \end{tikzpicture}
        \end{mlbox}
    \end{center}
\end{frame}

% ============================================================================
% AUFGABE 1: Werte ablesen
% ============================================================================

% --- Aufgabe 1: Frage ---
\begin{frame}{Werte aus dem Diagramm ablesen}
    \begin{columns}
        \begin{column}{0.45\textwidth}
            \textbf{Ein Radfahrer fährt gleichförmig.}

            \vspace{0.8em}
            \textbf{Lies ab:}
            \begin{enumerate}
                \item[a)] Weg nach $t = 4$ s?
                \item[b)] Zeit bei $s = 30$ m?
                \item[c)] Weg nach $t = 8$ s?
            \end{enumerate}
        \end{column}
        \begin{column}{0.5\textwidth}
            \begin{tikzpicture}[scale=0.5]
                \draw[gridcolor, very thin] (0,0) grid (8,5);
                \draw[->, thick] (0,0) -- (8.5,0) node[right] {$t$ in s};
                \draw[->, thick] (0,0) -- (0,5.5) node[above] {$s$ in m};
                \foreach \x in {0,2,4,6,8} \draw (\x,0) node[below] {\small \x};
                \foreach \y/\l in {0/0,1/10,2/20,3/30,4/40,5/50} \draw (0,\y) node[left] {\small \l};
                \draw[physik, very thick] (0,0) -- (8,4);
                \foreach \x/\y in {0/0,2/1,4/2,6/3,8/4} \filldraw[physik] (\x,\y) circle (3pt);
            \end{tikzpicture}
        \end{column}
    \end{columns}
\end{frame}

% --- Aufgabe 1: Lösung ---
\begin{frame}{Werte aus dem Diagramm ablesen}
    \begin{columns}
        \begin{column}{0.45\textwidth}
            \textbf{Ein Radfahrer fährt gleichförmig.}

            \vspace{0.8em}
            \textcolor{loesung}{\textbf{Lösung:}
            \begin{enumerate}
                \item[a)] $s = 20$ m
                \item[b)] $t = 6$ s
                \item[c)] $s = 40$ m
            \end{enumerate}}
        \end{column}
        \begin{column}{0.5\textwidth}
            \begin{tikzpicture}[scale=0.5]
                \draw[gridcolor, very thin] (0,0) grid (8,5);
                \draw[->, thick] (0,0) -- (8.5,0) node[right] {$t$ in s};
                \draw[->, thick] (0,0) -- (0,5.5) node[above] {$s$ in m};
                \foreach \x in {0,2,4,6,8} \draw (\x,0) node[below] {\small \x};
                \foreach \y/\l in {0/0,1/10,2/20,3/30,4/40,5/50} \draw (0,\y) node[left] {\small \l};
                \draw[physik, very thick] (0,0) -- (8,4);
                \draw[loesung, dashed] (4,0) -- (4,2) -- (0,2);
                \filldraw[loesung] (4,2) circle (4pt);
            \end{tikzpicture}
        \end{column}
    \end{columns}
\end{frame}

% --- Aufgabe 1: + AB-Verweis ---
\begin{frame}{Werte aus dem Diagramm ablesen}
    \begin{columns}
        \begin{column}{0.45\textwidth}
            \textbf{Ein Radfahrer fährt gleichförmig.}

            \vspace{0.8em}
            \textcolor{loesung}{\textbf{Lösung:}
            \begin{enumerate}
                \item[a)] $s = 20$ m
                \item[b)] $t = 6$ s
                \item[c)] $s = 40$ m
            \end{enumerate}}
        \end{column}
        \begin{column}{0.5\textwidth}
            \begin{tikzpicture}[scale=0.5]
                \draw[gridcolor, very thin] (0,0) grid (8,5);
                \draw[->, thick] (0,0) -- (8.5,0) node[right] {$t$ in s};
                \draw[->, thick] (0,0) -- (0,5.5) node[above] {$s$ in m};
                \foreach \x in {0,2,4,6,8} \draw (\x,0) node[below] {\small \x};
                \foreach \y/\l in {0/0,1/10,2/20,3/30,4/40,5/50} \draw (0,\y) node[left] {\small \l};
                \draw[physik, very thick] (0,0) -- (8,4);
                \draw[loesung, dashed] (4,0) -- (4,2) -- (0,2);
                \filldraw[loesung] (4,2) circle (4pt);
            \end{tikzpicture}
        \end{column}
    \end{columns}

    \vspace{0.3em}
    \begin{abhinweis}
        \textbf{Aufgabe 1:} Lies die Werte ab
    \end{abhinweis}
\end{frame}

% --- ML Aufgabe 1 ---
\begin{frame}{Musterlösung: Aufgabe 1}
    \begin{center}
        \begin{mlbox}
            \textbf{Aufgabe 1: Werte ablesen}

            \vspace{0.5em}
            a) Weg nach $t = 4$ s? \quad $s =$ \textcolor{loesung}{\textbf{20 m}}

            \vspace{0.3em}
            b) Zeit bei $s = 30$ m? \quad $t =$ \textcolor{loesung}{\textbf{6 s}}

            \vspace{0.3em}
            c) Weg nach $t = 8$ s? \quad $s =$ \textcolor{loesung}{\textbf{40 m}}
        \end{mlbox}
    \end{center}
\end{frame}

% ============================================================================
% AUFGABE 2: Diagramm zeichnen
% ============================================================================

% --- Aufgabe 2: Frage ---
\begin{frame}{Ein Diagramm selbst zeichnen}
    \begin{columns}
        \begin{column}{0.28\textwidth}
            \textbf{Ein Jogger läuft:}

            \vspace{0.3em}
            \begin{tabular}{|c|c|}
                \hline
                $t$ (s) & $s$ (m) \\
                \hline
                0 & 0 \\ 4 & 20 \\ 8 & 40 \\ 12 & 60 \\ 16 & 80 \\
                \hline
            \end{tabular}

            \vspace{0.5em}
            \textbf{So geht's:}
            \begin{enumerate}
                \item Punkte einzeichnen
                \item Mit Lineal verbinden
            \end{enumerate}
        \end{column}
        \begin{column}{0.68\textwidth}
            \begin{tikzpicture}[scale=0.42]
                \foreach \x in {0,1,...,16} \draw[gridcolor] (\x,0) -- (\x,8);
                \foreach \y in {0,1,...,8} \draw[gridcolor] (0,\y) -- (16,\y);
                \draw[->, thick] (0,0) -- (16.8,0) node[right] {$t$ in s};
                \draw[->, thick] (0,0) -- (0,8.8) node[above] {$s$ in m};
                \foreach \x/\l in {0/0,4/4,8/8,12/12,16/16} \node[below] at (\x,-0.3) {\small \l};
                \foreach \y/\l in {0/0,2/20,4/40,6/60,8/80} \node[left] at (-0.3,\y) {\small \l};
            \end{tikzpicture}
        \end{column}
    \end{columns}
\end{frame}

% --- Aufgabe 2: Lösung ---
\begin{frame}{Ein Diagramm selbst zeichnen}
    \begin{columns}
        \begin{column}{0.28\textwidth}
            \textbf{Ein Jogger läuft:}

            \vspace{0.3em}
            \begin{tabular}{|c|c|}
                \hline
                $t$ (s) & $s$ (m) \\
                \hline
                0 & 0 \\ 4 & 20 \\ 8 & 40 \\ 12 & 60 \\ 16 & 80 \\
                \hline
            \end{tabular}

            \vspace{0.5em}
            \textcolor{loesung}{\textbf{Ergebnis:}\\Eine Gerade durch den Ursprung!}
        \end{column}
        \begin{column}{0.68\textwidth}
            \begin{tikzpicture}[scale=0.42]
                \foreach \x in {0,1,...,16} \draw[gridcolor] (\x,0) -- (\x,8);
                \foreach \y in {0,1,...,8} \draw[gridcolor] (0,\y) -- (16,\y);
                \draw[->, thick] (0,0) -- (16.8,0) node[right] {$t$ in s};
                \draw[->, thick] (0,0) -- (0,8.8) node[above] {$s$ in m};
                \foreach \x/\l in {0/0,4/4,8/8,12/12,16/16} \node[below] at (\x,-0.3) {\small \l};
                \foreach \y/\l in {0/0,2/20,4/40,6/60,8/80} \node[left] at (-0.3,\y) {\small \l};
                \draw[loesung, very thick] (0,0) -- (16,8);
                \foreach \x/\y in {0/0,4/2,8/4,12/6,16/8} \filldraw[loesung] (\x,\y) circle (4pt);
            \end{tikzpicture}
        \end{column}
    \end{columns}
\end{frame}

% --- Aufgabe 2: + AB-Verweis ---
\begin{frame}{Ein Diagramm selbst zeichnen}
    \begin{columns}
        \begin{column}{0.28\textwidth}
            \textbf{Ein Jogger läuft:}

            \vspace{0.3em}
            \begin{tabular}{|c|c|}
                \hline
                $t$ (s) & $s$ (m) \\
                \hline
                0 & 0 \\ 4 & 20 \\ 8 & 40 \\ 12 & 60 \\ 16 & 80 \\
                \hline
            \end{tabular}

            \vspace{0.5em}
            \textcolor{loesung}{\textbf{Ergebnis:}\\Eine Gerade durch den Ursprung!}
        \end{column}
        \begin{column}{0.68\textwidth}
            \begin{tikzpicture}[scale=0.42]
                \foreach \x in {0,1,...,16} \draw[gridcolor] (\x,0) -- (\x,8);
                \foreach \y in {0,1,...,8} \draw[gridcolor] (0,\y) -- (16,\y);
                \draw[->, thick] (0,0) -- (16.8,0) node[right] {$t$ in s};
                \draw[->, thick] (0,0) -- (0,8.8) node[above] {$s$ in m};
                \foreach \x/\l in {0/0,4/4,8/8,12/12,16/16} \node[below] at (\x,-0.3) {\small \l};
                \foreach \y/\l in {0/0,2/20,4/40,6/60,8/80} \node[left] at (-0.3,\y) {\small \l};
                \draw[loesung, very thick] (0,0) -- (16,8);
                \foreach \x/\y in {0/0,4/2,8/4,12/6,16/8} \filldraw[loesung] (\x,\y) circle (4pt);
            \end{tikzpicture}
        \end{column}
    \end{columns}

    \vspace{0.3em}
    \begin{abhinweis}
        \textbf{Aufgabe 2:} Zeichne das Diagramm
    \end{abhinweis}
\end{frame}

% --- ML Aufgabe 2 ---
\begin{frame}{Musterlösung: Aufgabe 2}
    \begin{center}
        \begin{mlbox}
            \textbf{Aufgabe 2: Diagramm zeichnen}

            \vspace{0.5em}
            Punkte bei (0,0), (4,20), (8,40), (12,60), (16,80) eingetragen\\
            und mit dem Lineal zu einer \textcolor{loesung}{\textbf{Geraden}} verbunden.
        \end{mlbox}
    \end{center}
\end{frame}

% ============================================================================
% AUFGABE 3: Bewegungsabschnitte
% ============================================================================

% --- Aufgabe 3: Frage ---
\begin{frame}{Was passiert in den drei Abschnitten?}
    \begin{columns}
        \begin{column}{0.5\textwidth}
            \textbf{Beschreibe jeden Abschnitt:}

            \vspace{1em}
            Abschnitt 1: \luecke{3.5cm}

            \vspace{0.8em}
            Abschnitt 2: \luecke{3.5cm}

            \vspace{0.8em}
            Abschnitt 3: \luecke{3.5cm}

            \vspace{1em}
            \textbf{Nutze:} schnell, langsam, Stillstand
        \end{column}
        \begin{column}{0.45\textwidth}
            \begin{tikzpicture}[scale=0.55]
                \draw[gridcolor, very thin] (0,0) grid (8,5);
                \draw[->, thick] (0,0) -- (8.5,0) node[right] {$t$};
                \draw[->, thick] (0,0) -- (0,5.5) node[above] {$s$};
                \draw[physik, very thick] (0,0) -- (2,2.5);
                \draw[physik, very thick] (2,2.5) -- (5,2.5);
                \draw[physik, very thick] (5,2.5) -- (8,3.5);
                \node at (1,1.8) {\small 1};
                \node at (3.5,3) {\small 2};
                \node at (6.5,3.3) {\small 3};
            \end{tikzpicture}
        \end{column}
    \end{columns}
\end{frame}

% --- Aufgabe 3: Lösung ---
\begin{frame}{Was passiert in den drei Abschnitten?}
    \begin{columns}
        \begin{column}{0.5\textwidth}
            \textcolor{loesung}{\textbf{Lösung:}

            \vspace{0.5em}
            1: schnelle Bewegung (steil)

            \vspace{0.5em}
            2: Stillstand (waagerecht)

            \vspace{0.5em}
            3: langsame Bewegung (flach)}
        \end{column}
        \begin{column}{0.45\textwidth}
            \begin{tikzpicture}[scale=0.55]
                \draw[gridcolor, very thin] (0,0) grid (8,5);
                \draw[->, thick] (0,0) -- (8.5,0) node[right] {$t$};
                \draw[->, thick] (0,0) -- (0,5.5) node[above] {$s$};
                \draw[loesung, very thick] (0,0) -- (2,2.5);
                \draw[loesung, very thick] (2,2.5) -- (5,2.5);
                \draw[loesung, very thick] (5,2.5) -- (8,3.5);
                \node[loesung] at (1,1.5) {\tiny schnell};
                \node[loesung] at (3.5,2.9) {\tiny Stillst.};
                \node[loesung] at (6.5,3.3) {\tiny langs.};
            \end{tikzpicture}
        \end{column}
    \end{columns}
\end{frame}

% --- Aufgabe 3: + AB-Verweis ---
\begin{frame}{Was passiert in den drei Abschnitten?}
    \begin{columns}
        \begin{column}{0.5\textwidth}
            \textcolor{loesung}{\textbf{Lösung:}

            \vspace{0.5em}
            1: schnelle Bewegung (steil)

            \vspace{0.5em}
            2: Stillstand (waagerecht)

            \vspace{0.5em}
            3: langsame Bewegung (flach)}
        \end{column}
        \begin{column}{0.45\textwidth}
            \begin{tikzpicture}[scale=0.55]
                \draw[gridcolor, very thin] (0,0) grid (8,5);
                \draw[->, thick] (0,0) -- (8.5,0) node[right] {$t$};
                \draw[->, thick] (0,0) -- (0,5.5) node[above] {$s$};
                \draw[loesung, very thick] (0,0) -- (2,2.5);
                \draw[loesung, very thick] (2,2.5) -- (5,2.5);
                \draw[loesung, very thick] (5,2.5) -- (8,3.5);
                \node[loesung] at (1,1.5) {\tiny schnell};
                \node[loesung] at (3.5,2.9) {\tiny Stillst.};
                \node[loesung] at (6.5,3.3) {\tiny langs.};
            \end{tikzpicture}
        \end{column}
    \end{columns}

    \vspace{0.3em}
    \begin{abhinweis}
        \textbf{Aufgabe 3:} Beschreibe die drei Abschnitte
    \end{abhinweis}
\end{frame}

% --- ML Aufgabe 3 ---
\begin{frame}{Musterlösung: Aufgabe 3}
    \begin{center}
        \begin{mlbox}
            \textbf{Aufgabe 3: Bewegungsabschnitte}

            \vspace{0.5em}
            1: \textcolor{loesung}{\textbf{schnelle gleichförmige Bewegung}} (steil)

            \vspace{0.3em}
            2: \textcolor{loesung}{\textbf{Stillstand}} (waagerecht)

            \vspace{0.3em}
            3: \textcolor{loesung}{\textbf{langsame gleichförmige Bewegung}} (flach)
        \end{mlbox}
    \end{center}
\end{frame}

% ============================================================================
% AUFGABE 4: Stillstand oder gleichförmig?
% ============================================================================

% --- Aufgabe 4: Frage ---
\begin{frame}{Stillstand oder gleichförmig?}
    \begin{columns}
        \begin{column}{0.5\textwidth}
            \textbf{Welche Bewegungsart zeigen A und B?}

            \vspace{1.5em}
            Linie A:\\
            \hspace{1em} $\square$ gleichförmig \quad $\square$ Stillstand

            \vspace{1em}
            Linie B:\\
            \hspace{1em} $\square$ gleichförmig \quad $\square$ Stillstand

            \vspace{1em}
            \textbf{Tipp:} Wo liegt Linie A?
        \end{column}
        \begin{column}{0.45\textwidth}
            \begin{tikzpicture}[scale=0.55]
                \draw[gridcolor, very thin] (0,0) grid (8,5);
                \draw[->, thick] (0,0) -- (8.5,0) node[right] {$t$};
                \draw[->, thick] (0,0) -- (0,5.5) node[above] {$s$};
                \draw[physik, very thick] (0,0) -- (8,0);
                \node[physik, right] at (8.2,0) {\small A};
                \draw[red!70!black, very thick] (0,0) -- (8,3.5);
                \node[red!70!black, right] at (8.2,3.5) {\small B};
            \end{tikzpicture}
        \end{column}
    \end{columns}
\end{frame}

% --- Aufgabe 4: Lösung ---
\begin{frame}{Stillstand oder gleichförmig?}
    \begin{columns}
        \begin{column}{0.5\textwidth}
            \textcolor{loesung}{\textbf{Lösung:}}

            \vspace{1em}
            \textcolor{loesung}{A: \textbf{Stillstand}\\
            (liegt auf x-Achse, $s = 0$)}

            \vspace{1em}
            \textcolor{loesung}{B: \textbf{gleichförmig}\\
            (Gerade mit Steigung)}
        \end{column}
        \begin{column}{0.45\textwidth}
            \begin{tikzpicture}[scale=0.55]
                \draw[gridcolor, very thin] (0,0) grid (8,5);
                \draw[->, thick] (0,0) -- (8.5,0) node[right] {$t$};
                \draw[->, thick] (0,0) -- (0,5.5) node[above] {$s$};
                \draw[loesung, very thick] (0,0) -- (8,0);
                \node[loesung, above] at (4,0.3) {\tiny Stillstand};
                \draw[loesung, very thick] (0,0) -- (8,3.5);
                \node[loesung, above] at (6,3) {\tiny gleichf.};
            \end{tikzpicture}
        \end{column}
    \end{columns}
\end{frame}

% --- Aufgabe 4: + AB-Verweis ---
\begin{frame}{Stillstand oder gleichförmig?}
    \begin{columns}
        \begin{column}{0.5\textwidth}
            \textcolor{loesung}{\textbf{Lösung:}}

            \vspace{1em}
            \textcolor{loesung}{A: \textbf{Stillstand}\\
            (liegt auf x-Achse, $s = 0$)}

            \vspace{1em}
            \textcolor{loesung}{B: \textbf{gleichförmig}\\
            (Gerade mit Steigung)}
        \end{column}
        \begin{column}{0.45\textwidth}
            \begin{tikzpicture}[scale=0.55]
                \draw[gridcolor, very thin] (0,0) grid (8,5);
                \draw[->, thick] (0,0) -- (8.5,0) node[right] {$t$};
                \draw[->, thick] (0,0) -- (0,5.5) node[above] {$s$};
                \draw[loesung, very thick] (0,0) -- (8,0);
                \node[loesung, above] at (4,0.3) {\tiny Stillstand};
                \draw[loesung, very thick] (0,0) -- (8,3.5);
                \node[loesung, above] at (6,3) {\tiny gleichf.};
            \end{tikzpicture}
        \end{column}
    \end{columns}

    \vspace{0.3em}
    \begin{abhinweis}
        \textbf{Aufgabe 4:} Kreuze an
    \end{abhinweis}
\end{frame}

% --- ML Aufgabe 4 ---
\begin{frame}{Musterlösung: Aufgabe 4}
    \begin{center}
        \begin{mlbox}
            \textbf{Aufgabe 4: Stillstand oder gleichförmig?}

            \vspace{0.5em}
            A: $\square$ gleichförmig \quad \textcolor{loesung}{$\boxtimes$ \textbf{Stillstand}}

            \vspace{0.3em}
            B: \textcolor{loesung}{$\boxtimes$ \textbf{gleichförmig}} \quad $\square$ Stillstand
        \end{mlbox}
    \end{center}
\end{frame}

% ============================================================================
% ZUSAMMENFASSUNG
% ============================================================================
\begin{frame}{Zusammenfassung}
    \begin{columns}
        \begin{column}{0.55\textwidth}
            \textbf{Das Weg-Zeit-Diagramm:}
            \begin{itemize}
                \item Zeit $t$ $\rightarrow$ waagerechte Achse
                \item Weg $s$ $\rightarrow$ senkrechte Achse
                \item Gleichförmig $\rightarrow$ Gerade
            \end{itemize}

            \vspace{0.5em}
            \textbf{Die Steigung zeigt die Geschwindigkeit:}
            \begin{itemize}
                \item Steil = schnell
                \item Flach = langsam
                \item Waagerecht = Stillstand
                \item Kurve = ungleichförmig
            \end{itemize}
        \end{column}
        \begin{column}{0.4\textwidth}
            \begin{tikzpicture}[scale=0.38]
                \draw[gridcolor, very thin] (0,0) grid (8,6);
                \draw[->, thick] (0,0) -- (8.5,0) node[right] {$t$};
                \draw[->, thick] (0,0) -- (0,6.5) node[above] {$s$};
                \draw[red!70!black, very thick] (0,0) -- (3,5);
                \node[red!70!black, right] at (3,5) {\tiny schnell};
                \draw[physik, very thick] (0,0) -- (8,2.5);
                \node[physik, right] at (8,2.5) {\tiny langsam};
                \draw[gray, very thick] (0,4) -- (8,4);
                \node[gray, right] at (8,4) {\tiny Stillstand};
            \end{tikzpicture}
        \end{column}
    \end{columns}

    \vspace{0.5em}
    \begin{abhinweis}
        Prüfe: Sind alle Kästen und Aufgaben vollständig?
    \end{abhinweis}
\end{frame}

\end{document}
